\section{Top-level AXI4 interface}
\label{sec:axi4top}

CIF presents a standard AXI4 slave interface to its client(s).  The table
below cross-references the channel signals against the widths the current
architecture documents establish (\texttt{ADDR\_W=32}, \texttt{AXID\_W=6},
\texttt{LEN\_W=8}, from the Request Validator's Design Parameters).  Widths
marked AXI4-fixed follow the AXI4 protocol specification itself and are not
CIF-specific parameters.  Write/read data-path widths are not established by
any current CIF document and remain \textbf{TBD}; see \S\ref{sec:respformat}
for the corresponding response-side discussion.

\begin{longtable}{p{1.3cm}p{3.2cm}p{1.6cm}p{9.5cm}}
\toprule
\textbf{Channel} & \textbf{Signal} & \textbf{Width} & \textbf{Notes}\\
\midrule
AW & \texttt{AWVALID/AWREADY} & 1/1 & Write address handshake\\
   & \texttt{AWADDR}          & 32  & \texttt{ADDR\_W}\\
   & \texttt{AWID}            & 6   & \texttt{AXID\_W}\\
   & \texttt{AWLEN}           & 8   & \texttt{LEN\_W}; \emph{not} the MC packet/segment count (\S\ref{sec:identity})\\
   & \texttt{AWSIZE/AWBURST}  & 3/2 & AXI4-fixed\\
\midrule
W  & \texttt{WVALID/WREADY}   & 1/1 & Write data handshake\\
   & \texttt{WDATA/WSTRB}     & TBD & Data-path width not established by current CIF documents (\S\ref{sec:pktformat})\\
   & \texttt{WLAST}           & 1   & AXI4-fixed\\
\midrule
B  & \texttt{BVALID/BREADY}   & 1/1 & Write response handshake\\
   & \texttt{BID}             & 6   & \texttt{AXID\_W}; from the retiring transaction's Slot FIFO context, not stored in the Metadata RAM entry\\
   & \texttt{BRESP}           & 2   & AXI4-fixed; driven from the write ROB's aggregate \texttt{resp\_status}\\
\midrule
AR & \texttt{ARVALID/ARREADY} & 1/1 & Read address handshake\\
   & \texttt{ARADDR}          & 32  & \texttt{ADDR\_W}\\
   & \texttt{ARID}            & 6   & \texttt{AXID\_W}\\
   & \texttt{ARLEN}           & 8   & \texttt{LEN\_W}; \emph{not} the MC packet/segment count (\S\ref{sec:identity})\\
   & \texttt{ARSIZE/ARBURST}  & 3/2 & AXI4-fixed\\
\midrule
R  & \texttt{RVALID/RREADY}   & 1/1 & Read data handshake\\
   & \texttt{RDATA}           & TBD & Data-path width not established by current CIF documents (\S\ref{sec:pktformat})\\
   & \texttt{RID}             & 6   & \texttt{AXID\_W}; from the retiring transaction's Slot FIFO context\\
   & \texttt{RRESP}           & 2   & AXI4-fixed; per-beat value from the read ROB's serialized status\\
   & \texttt{RLAST}           & 1   & Asserted only on the final AXI beat of the final segment, not on every segment boundary\\
\bottomrule
\caption{Top-level AXI4 slave interface, cross-referenced to CIF-internal handling}
\end{longtable}

\section{Interface Signal Table}
\label{sec:interfaces}

This table records the consolidated interfaces after the read/write ROB
refactor.  Names describe the frozen connectivity where prior documents did
not fix a signal name.

\subsection{AXI request acceptance to the selected ROB}
\begin{longtable}{p{3.3cm}p{1cm}p{3.4cm}p{7cm}}
\toprule
\textbf{Signal} & \textbf{Width} & \textbf{Path} & \textbf{Description}\\
\midrule
\texttt{alloc\_req} & 1 & Accept logic $\rightarrow$ RD/WR ROB & Allocate one transaction slot on accepted AR/AW request.\\
\texttt{AXID} & 6 & Accept logic $\rightarrow$ RD/WR ROB & Selects the per-AXID Slot FIFO.\\
\texttt{ROB\_INDEX} & 8 & RD/WR ROB $\rightarrow$ Splitter & Free-List result retained for the transaction.\\
\texttt{splitter\_pkt\_count} & 3 & Splitter $\rightarrow$ RD/WR ROB & Number of segments for this transaction, 1--4; initializes \texttt{exp\_pkts}.\\
\bottomrule
\end{longtable}

\subsection{Burst Splitter to Packet Formation}
\begin{longtable}{p{3.3cm}p{1cm}p{3.4cm}p{7cm}}
\toprule
\textbf{Signal} & \textbf{Width} & \textbf{Path} & \textbf{Description}\\
\midrule
\texttt{seg\_valid/seg\_ready} & 1/1 & Splitter $\leftrightarrow$ Packet Formation & Handshake for the current Segment Table entry.\\
\texttt{seg\_addr} & 32 & Splitter $\rightarrow$ Packet Formation & Segment start address.\\
\texttt{seg\_beat\_cnt} & 8 & Splitter $\rightarrow$ Packet Formation & Number of AXI beats represented by the segment.\\
\texttt{PACKET\_NUMBER} & 2 & Splitter $\rightarrow$ Packet Formation & Current \texttt{seg\_idx}, assigned at packet emission.\\
\texttt{GLOBAL\_TAG} & 10 & Packet Formation $\rightarrow$ Packet FIFO & \texttt{\{ROB\_INDEX,PACKET\_NUMBER\}}.\\
\bottomrule
\end{longtable}

\subsection{MC Core completion interface}

Response transport and handshake naming are MC/Data-Buffer
interface-contract dependent (\S\ref{sec:respformat}); no
\texttt{RESP\_VALID}/\texttt{RESP\_READY} signal pair is assumed.  The fields
below are the only ones the current architecture requires MC Core to supply.

\begin{longtable}{p{3.3cm}p{1cm}p{3.4cm}p{7cm}}
\toprule
\textbf{Signal} & \textbf{Width} & \textbf{Path} & \textbf{Description}\\
\midrule
\texttt{GLOBAL\_TAG} & 10 & MC Core $\rightarrow$ RD/WR ROB & Echoed request tag; directly decodes index and packet number.\\
\texttt{data} & TBD & MC Core $\rightarrow$ Read ROB & Read data payload; width and response-unit granularity are contract-dependent (\S\ref{sec:respformat}).\\
\texttt{status} & 4 & MC Core $\rightarrow$ RD/WR ROB & Raw completion status; encoding is \textbf{TBD with MC Core}.\\
\bottomrule
\end{longtable}

\subsection{Retirement interfaces}
\begin{longtable}{p{3.3cm}p{1cm}p{3.4cm}p{7cm}}
\toprule
\textbf{Signal} & \textbf{Width} & \textbf{Path} & \textbf{Description}\\
\midrule
\texttt{merge\_ready} & 1 & Merge Logic $\rightarrow$ Read ROB & Accepts an emitted AXI read beat; serializer progress and Data Buffer release follow the MC/Data-Buffer contract.\\
\texttt{rob\_last} & 1 & Read ROB $\rightarrow$ Merge Logic & Asserted with the final serialized read beat.\\
\texttt{bresp\_ready} & 1 & B response path $\rightarrow$ Write ROB & Accepts the ordered BRESP and permits write-slot release.\\
\texttt{dbuf\_alloc/read/free} & TBD & Read ROB $\leftrightarrow$ Read Data Buffer & Allocation representation, handshake, response-unit semantics, and release timing are \textbf{TBD}; one handle is recorded per packet number.\\
\bottomrule
\end{longtable}

\section{Arbitration and backpressure}

The request validators and Burst Splitter retain the existing request
arbitration policy.  The selected ROB Free List and selected AXID Slot FIFO
provide transaction-allocation backpressure.  A full Slot FIFO affects only
its AXID; it does not imply global Free List exhaustion.  Conversely, an empty
Free List means all 256 transaction entries are allocated.

Read packet issuance is additionally controlled by pooled Read Data Buffer
availability and its independent 128-entry watermark.  Metadata RAM capacity
does not control packet issue.

\section{Reset and initialization sequence}
\label{sec:resetseq}

Reset is active-low, asserted asynchronously, and deasserted synchronously to
the CIF clock domain.  While asserted, it forces every structure below to its
initialized state; normal operation begins only after synchronous
deassertion.

On reset, each ROB instance independently initializes its 256-bit Free List
to free, clears Metadata RAM state, resets all 64 Slot FIFOs, and clears its
Retiring Register.  The Burst Splitter clears its Segment Table validity,
\texttt{num\_segments}, \texttt{seg\_idx}, and retained \texttt{ROB\_INDEX}.
Packet FIFOs and AXI request FIFOs retain their established reset behavior.
The MC completion/response handshake is not assigned a fixed reset value by
this specification; its behavior across reset is governed by the MC--CIF
interface contract (\S\ref{sec:respformat}).

A soft reset first quiesces request acceptance and waits for packet FIFOs,
Segment Table activity, both ROB Slot-FIFO sets, both Retiring Registers, and
all externally visible AXI responses to drain.  The Read Data Buffer has a
separate drain/initialization contract; its exact interface and timing are
\textbf{TBD}.  Reset must not release a read \texttt{ROB\_INDEX} until the
final AXI read response and its Data Buffer release have completed.

\section{Error propagation}

Error packets are segmented and tagged in the same manner as ordinary
requests.  Every completion sets its packet-number bit in the selected ROB.
The write ROB aggregates packet status before its one BRESP; the read ROB
serializes status with its data.  Raw MC status encoding and write
worst-status precedence are \textbf{TBD}; this specification intentionally
does not assign an AXI response value without that contract.
